\documentclass[10pt]{article}
% Shared LaTeX style for MIT 18.06 exam revision notes.
% Compile topic files with Tectonic, XeLaTeX, or LuaLaTeX.

\usepackage[a4paper,margin=0.72in]{geometry}
\usepackage{fontspec}
\usepackage{amsmath,amssymb,mathtools}
\usepackage{booktabs,array,tabularx}
\usepackage{enumitem}
\usepackage{titlesec}
\usepackage{fancyhdr}
\usepackage{microtype}
\usepackage{xcolor}

\setmainfont{Times New Roman}
\setsansfont{Arial}

\setlength{\parindent}{0pt}
\setlength{\parskip}{3.6pt}
\setlength{\abovedisplayskip}{6pt}
\setlength{\belowdisplayskip}{6pt}
\setlength{\abovedisplayshortskip}{4pt}
\setlength{\belowdisplayshortskip}{4pt}
\renewcommand{\arraystretch}{1.08}
\setlength{\tabcolsep}{4.5pt}

\titleformat{\section}
  {\normalfont\large\bfseries}
  {\thesection}
  {0.55em}
  {}
\titlespacing*{\section}{0pt}{10pt}{3pt}

\titleformat{\subsection}
  {\normalfont\normalsize\bfseries}
  {\thesubsection}
  {0.5em}
  {}
\titlespacing*{\subsection}{0pt}{7pt}{2pt}

\setlist[itemize]{leftmargin=1.35em,itemsep=1pt,topsep=2pt,parsep=0pt}
\setlist[enumerate]{leftmargin=1.55em,itemsep=1pt,topsep=2pt,parsep=0pt}

\pagestyle{fancy}
\fancyhf{}
\fancyfoot[C]{\scriptsize\color{black!45}MIT 18.06 Linear Algebra Exam Notes --- \thepage}
\renewcommand{\headrulewidth}{0pt}
\renewcommand{\footrulewidth}{0pt}

\newcommand{\notesubtitle}[1]{%
  {\centering\itshape #1\par}
  \vspace{2pt}
}

\newcommand{\source}[1]{%
  {\centering\small #1\par}
  \vspace{6pt}
}

\newcommand{\labelnote}[2]{%
  \par\smallskip
  \noindent\textbf{#1.}\quad #2
  \par\smallskip
}

\newcommand{\tightlabel}[1]{%
  \par\smallskip
  \noindent\textbf{#1.}\quad
}

\newcolumntype{Y}{>{\raggedright\arraybackslash}X}
\newcolumntype{L}[1]{>{\raggedright\arraybackslash}p{#1}}

\newenvironment{notetable}
  {\small\begin{tabularx}{\linewidth}}
  {\end{tabularx}}


\begin{document}

\begin{center}
{\LARGE 6.1--6.2 Eigenvalues, Eigenvectors, and Diagonalization\par}
\end{center}
\notesubtitle{Concise exam revision notes for eigenvectors, characteristic equations, powers, and stability.}
\source{Based on handwritten MIT 18.06 revision notes.}

\section{Core Idea}

\labelnote{Core idea}{Most vectors change direction when multiplied by \(A\). Eigenvectors keep their direction; only their length and sign are scaled.}

The defining equation is
\[
Ax=\lambda x,
\qquad x\ne0.
\]
Equivalently,
\[
(A-\lambda I)x=0.
\]
So \(A-\lambda I\) must be singular:
\[
\det(A-\lambda I)=0.
\]

\section{How to Find Eigenvalues and Eigenvectors}

\begin{enumerate}
  \item Solve \(\det(A-\lambda I)=0\) to find eigenvalues.
  \item For each eigenvalue, solve \((A-\lambda I)x=0\).
  \item Use any nonzero vector in that nullspace as an eigenvector.
\end{enumerate}

\labelnote{Exam trigger}{Eigenvectors come from a nullspace. If the nullspace has more than one direction, choose a convenient basis.}

\section{Fast Checks}

\begin{center}
\small
\begin{tabularx}{\linewidth}{@{}L{0.25\linewidth}Y Y@{}}
\toprule
Check & Formula & Meaning \\
\midrule
Trace & \(\operatorname{trace}(A)=\lambda_1+\cdots+\lambda_n\) & Sum of eigenvalues. \\
Determinant & \(\det A=\lambda_1\lambda_2\cdots\lambda_n\) & Product of eigenvalues. \\
Invertibility & \(0\) is not an eigenvalue & Equivalent to \(\det A\ne0\). \\
\bottomrule
\end{tabularx}
\end{center}

These checks catch arithmetic mistakes after solving the characteristic equation.

\section{Special Matrices}

\begin{center}
\small
\begin{tabularx}{\linewidth}{@{}L{0.26\linewidth}Y Y@{}}
\toprule
Matrix type & Typical eigenvalues & Exam memory \\
\midrule
Projection \(P\) & \(0\) and \(1\) & \(P^{2}=P\). \\
Reflection \(R\) & \(1\) and \(-1\) & Directions along and across the mirror. \\
Rotation in the plane & Usually complex & Real rotations may have no real eigenvectors. \\
Triangular matrix & Diagonal entries & Read eigenvalues from the diagonal. \\
\bottomrule
\end{tabularx}
\end{center}

\labelnote{Common mistake}{Row reduction changes eigenvalues in general. Do not row-reduce \(A\) and read eigenvalues from the new matrix.}

\section{Diagonalization}

If \(A\) has \(n\) independent eigenvectors, put them into
\[
S=\begin{bmatrix}x_1&x_2&\cdots&x_n\end{bmatrix}.
\]
Then
\[
AS=S\Lambda,
\qquad
A=S\Lambda S^{-1}.
\]
Here
\[
\Lambda=
\begin{bmatrix}
\lambda_1&0&\cdots&0\\
0&\lambda_2&\cdots&0\\
\vdots&\vdots&\ddots&\vdots\\
0&0&\cdots&\lambda_n
\end{bmatrix}.
\]

\begin{center}
\small
\begin{tabularx}{\linewidth}{@{}L{0.32\linewidth}Y@{}}
\toprule
Case & Diagonalization result \\
\midrule
\(n\) distinct eigenvalues & Automatically enough independent eigenvectors. \\
Repeated eigenvalue & May or may not have enough eigenvectors. \\
Not enough eigenvectors & Not diagonalizable. \\
\bottomrule
\end{tabularx}
\end{center}

\labelnote{Exam memory}{Repeated eigenvalues are not the problem by themselves. The issue is whether there are enough independent eigenvectors.}

\section{Powers and Difference Equations}

Diagonalization makes powers easy:
\[
A^{k}=S\Lambda^{k}S^{-1}.
\]
For a difference equation
\[
u_{k+1}=Au_k,
\]
the solution is
\[
u_k=A^ku_0=S\Lambda^kS^{-1}u_0.
\]
If
\[
u_0=c_1x_1+\cdots+c_nx_n,
\]
then
\[
u_k=c_1\lambda_1^kx_1+\cdots+c_n\lambda_n^kx_n.
\]

\section{Stability}

\begin{center}
\small
\begin{tabularx}{\linewidth}{@{}L{0.30\linewidth}Y@{}}
\toprule
Eigenvalue behavior & Effect on \(A^k\) \\
\midrule
All \(|\lambda_i|<1\) & Decays to zero. \\
Some \(|\lambda_i|>1\) & Blows up in that direction. \\
\(|\lambda_i|=1\) & Borderline; behavior depends on the matrix. \\
One \(\lambda=1\), rest \(|\lambda|<1\) & Possible steady state. \\
\bottomrule
\end{tabularx}
\end{center}

\section{Common Mistakes}

\begin{center}
\small
\begin{tabularx}{\linewidth}{@{}L{0.40\linewidth}Y@{}}
\toprule
Mistake & Correct exam habit \\
\midrule
Solving \(Ax=0\) for every eigenvector. & Solve \((A-\lambda I)x=0\). \\
Forgetting \(x\ne0\). & Eigenvectors are never the zero vector. \\
Assuming \(\lambda=0\) means diagonalizable fails. & It means \(A\) is singular; diagonalizability is separate. \\
Assuming repeated eigenvalues always fail. & Count independent eigenvectors. \\
Row-reducing \(A\) before finding eigenvalues. & Use \(\det(A-\lambda I)=0\). \\
Confusing invertibility and diagonalizability. & Invertible means no zero eigenvalue; diagonalizable means enough eigenvectors. \\
\bottomrule
\end{tabularx}
\end{center}

\section{Final Exam Checklist}

\tightlabel{Checklist}
\begin{enumerate}
  \item Compute \(\det(A-\lambda I)\).
  \item Solve for all eigenvalues \(\lambda\).
  \item Check trace and determinant against the eigenvalues.
  \item For each \(\lambda\), solve \((A-\lambda I)x=0\).
  \item Count independent eigenvectors.
  \item Diagonalize only if there are \(n\) independent eigenvectors.
  \item Form \(S\) from eigenvectors and \(\Lambda\) from eigenvalues.
  \item Use \(A^{k}=S\Lambda^{k}S^{-1}\) for powers.
  \item Use \(|\lambda|\) to decide growth or decay in difference equations.
\end{enumerate}

\end{document}
